\expandafter\ifx\csname CAPETINCLUDEMODE\endcsname\relax
\documentclass[11pt,a4paper]{article}
% preamble.tex — packages et commandes communes a toutes les fiches CAPET SII IE
\input{glyphtounicode}
\pdfgentounicode=1
\usepackage{cmap}
\usepackage[T1]{fontenc}
\usepackage[utf8]{inputenc}
\usepackage[french]{babel}
\usepackage[a4paper, top=1.6cm, bottom=1.6cm, left=1.4cm, right=1.4cm]{geometry}
\setlength{\headheight}{14.5pt}
\usepackage{amsmath, amssymb, amsfonts}
\usepackage{siunitx}
\sisetup{locale=FR, output-decimal-marker={,}, inter-unit-product=\ensuremath{{}\cdot{}}}
\usepackage{array, booktabs, multirow, tabularx}
\usepackage{xcolor}
\usepackage[most, breakable]{tcolorbox}
\usepackage{multicol}
\usepackage{enumitem}
\usepackage{mathtools}
\usepackage{lmodern}
\usepackage{microtype}
\usepackage{fancyhdr}
\usepackage{tikz}
\usetikzlibrary{positioning,arrows.meta}
\usepackage{lastpage}

% ── Couleurs par domaine ──────────────────────────────────────────────
\definecolor{colElec}{RGB}{26, 82, 148}       % Électrotechnique — bleu foncé
\definecolor{colElecLight}{RGB}{210, 225, 245} % fond bleu clair
\definecolor{colAuto}{RGB}{34, 120, 66}        % Automatique — vert
\definecolor{colMeca}{RGB}{185, 90, 20}        % Mécanique — orange
\definecolor{colTherm}{RGB}{160, 40, 40}       % Thermique — rouge
\definecolor{colQT}{RGB}{232, 244, 233}        % fond question type (vert très léger)
\definecolor{colJury}{RGB}{255, 243, 205}      % fond jury (jaune ambre)
\definecolor{colJuryFrame}{RGB}{200, 130, 0}   % bordure jury
\definecolor{colGrayLight}{RGB}{245, 245, 247} % fond gris neutre
\definecolor{colSeparator}{RGB}{180, 180, 180} % gris séparateur

% ── Style tcolorbox de section principale ────────────────────────────
\tcbset{
  ficheSection/.style={
    breakable,
    enhanced,
    colback=#1!7!white,
    colframe=#1,
    colbacktitle=#1,
    coltitle=white,
    fonttitle=\bfseries\normalsize,
    left=4pt, right=4pt, top=3pt, bottom=4pt,
    before skip=6pt, after skip=6pt,
  }
}

% ── Style encadré question type ──────────────────────────────────────
\tcbset{
  questiontype/.style={
    breakable,
    enhanced,
    colback=colQT,
    colframe=colElec!55!black,
    colbacktitle=colElec!55!black,
    coltitle=white,
    fonttitle=\bfseries\small,
    left=4pt, right=4pt, top=3pt, bottom=3pt,
    before skip=5pt, after skip=5pt,
  }
}

% ── Style encadré jury ───────────────────────────────────────────────
\tcbset{
  jury/.style={
    breakable,
    enhanced,
    colback=colJury,
    colframe=colJuryFrame,
    fonttitle=\bfseries\normalsize,
    left=4pt, right=4pt, top=3pt, bottom=3pt,
  }
}

% ── Mise en page ─────────────────────────────────────────────────────
\pagestyle{fancy}
\fancyhf{}
\renewcommand{\headrulewidth}{0.4pt}
\renewcommand{\footrulewidth}{0.3pt}

\setlist[itemize]{leftmargin=1.2em, itemsep=1pt, topsep=2pt, parsep=0pt}
\setlist[enumerate]{leftmargin=1.2em, itemsep=3pt, topsep=2pt, parsep=0pt}

\setlength{\parskip}{2pt}
\setlength{\columnsep}{10pt}

\def\CAPETBEGINFILE{\begin{document}}
\def\CAPETENDFILE{\end{document}}
\else
\def\CAPETBEGINFILE{}
\def\CAPETENDFILE{}
\fi

\fancyhead[L]{\textbf{\textcolor{colElec}{T06}} - Transformateur (monophase et triphase)}
\fancyhead[R]{\small\textit{CAPET SII IE - Session 2026}}
\fancyfoot[C]{\small\thepage\,/\,\pageref{LastPage}}

\CAPETBEGINFILE

% -----------------------------------------------------------------------
%  BLOC EN-TeTE
% -----------------------------------------------------------------------

\begin{tcolorbox}[
  enhanced, breakable,
  colback=colElec,
  colframe=colElec,
  left=6pt, right=6pt, top=5pt, bottom=5pt,
  before skip=0pt, after skip=4pt
]
  {\color{white}\bfseries\large T06 \quad Transformateur (monophase et triphase)}\\[2pt]
  {\color{colElec!30!white}\small
    \textbf{Domaine :} electrotechnique \quad
    \textbf{Presence :} \textbf{10/13} sujets (2013-2026) - theme \emph{tres frequent}
  }\\[1pt]
  {\color{colElec!20!white}\footnotesize
    Sous-themes : rapport de transformation - essai - vide ($R_f$, $X_m$) - essai court-circuit ($R_s$, $\varepsilon_{cc}$) - couplages Yy/Yd - chute de tension - rendement
  }
\end{tcolorbox}

% -----------------------------------------------------------------------
%  1. RAPPELS DE COURS
% -----------------------------------------------------------------------

\begin{tcolorbox}[
  ficheSection=colElec,
  title={1. Rappels de cours},
]

% -- 1.1 Rapport de transformation ------------------------------------
\textbf{1.1 Rapport de transformation}

\smallskip
\begin{center}
\renewcommand{\arraystretch}{1.45}
\begin{tabular}{lll}
\toprule
\textbf{Grandeur} & \textbf{Expression} & \textbf{Remarque} \\
\midrule
Rapport de transformation & $m = \dfrac{N_2}{N_1} = \dfrac{V_2}{V_1} = \dfrac{I_1}{I_2}$ & Monophase (tensions - vide) \\[8pt]
Rapport (triphase) & $m = \dfrac{U_{2n}}{U_{1n}}$ & Tensions composees etoile-toile \\[8pt]
Tension secondaire - vide & $U_s = m \cdot U_1$ & Sans charge \\[4pt]
Couplage Yy & $U_{composee} = \sqrt{3}\cdot U_{simple}$ & Rapport $m$ inchange \\[4pt]
Couplage Yd & Decalage de $30^\circ$ prim./sec. & elimine harmoniques rang 3 \\
\bottomrule
\end{tabular}
\end{center}

\medskip
\hrule\vspace{4pt}

% -- 1.2 Schema equivalent - essai - vide -----------------------------
\textbf{1.2 Essai - vide - pertes fer}

\smallskip - vide ($I_2 = 0$) : le schema equivalent se reduit - $R_f$ et $X_m$ en parallele, alimentes par $V_1$.
\[
  P_0 = \frac{V_1^2}{R_f} \quad\Rightarrow\quad R_f = \frac{V_1^2}{P_0}
  \qquad;\qquad
  Q_0 = \frac{V_1^2}{X_m} \quad\Rightarrow\quad X_m = \frac{V_1^2}{Q_0}
\]
\[
  S_0 = V_1 \cdot I_0 \qquad;\qquad Q_0 = \sqrt{S_0^2 - P_0^2}
\]
\textit{$P_0$ = pertes fer (mesurees au wattmetre cete primaire). $Q_0$ est toujours positif (charge inductive).}

\medskip
\hrule\vspace{4pt}

% -- 1.3 Essai en court-circuit - pertes cuivre -----------------------
\textbf{1.3 Essai en court-circuit - pertes cuivre}

\smallskip
\[
  P_{cc} = R_s \cdot I_{cc}^2 \quad\Rightarrow\quad R_s = R_1 + m^2 R_2 \quad \text{(resistances ramenees au primaire)}
\]
\[
  \varepsilon_{cc} = \frac{V_{cc}}{V_{1n}} \quad (\text{valeur typique : } 4\,\% \text{ - } 10\,\%)
\]
Chute de tension relative en charge :
\[
  \Delta U \approx \varepsilon_R \cos\varphi + \varepsilon_X \sin\varphi \qquad \text{avec}\quad
  \varepsilon_R = \frac{P_{cc}}{S_n},\quad \varepsilon_X \approx \varepsilon_{cc}
\]

\medskip
\hrule\vspace{4pt}

% -- 1.4 Rendement ----------------------------------------------------
\textbf{1.4 Rendement}

\smallskip
\[
  \eta = \frac{P_2}{P_1} = \frac{P_2}{P_2 + P_{fer} + P_{cuivre}}
\]
Maximum lorsque $P_{fer} = P_{cuivre}$ (fonctionnement au tiers de charge environ pour les transformateurs de distribution).

\end{tcolorbox}

% -----------------------------------------------------------------------
%  2. FORMULAIRE SYNTHeTIQUE
% -----------------------------------------------------------------------

\begin{tcolorbox}[
  ficheSection=colElec,
  title={2. Formulaire synthetique},
]
\small
\begin{multicols}{2}
\renewcommand{\arraystretch}{1.45}
\begin{tabular}{@{}p{4.5cm}ll@{}}
\toprule
\textbf{Formule} & \textbf{Usage} & \textbf{U.} \\
\midrule
\multicolumn{3}{@{}l}{\footnotesize\textit{Rapport de transformation}} \\
$m = N_2/N_1 = U_2/U_1 = I_1/I_2$ & Rapport & - \\
$U_s = m \cdot U_1$ & Tension sec. & V \\
\midrule
\multicolumn{3}{@{}l}{\footnotesize\textit{Essai - vide}} \\
$S_0 = V_1 I_0$ & Puiss. app. & VA \\
$Q_0 = \sqrt{S_0^2 - P_0^2}$ & Puiss. react. & var \\
$R_f = V_1^2 / P_0$ & Res. perte fer & $\Omega$ \\
$X_m = V_1^2 / Q_0$ & React. mag. & $\Omega$ \\
\bottomrule
\end{tabular}

\columnbreak

\renewcommand{\arraystretch}{1.45}
\begin{tabular}{@{}p{4.5cm}ll@{}}
\toprule
\textbf{Formule} & \textbf{Usage} & \textbf{U.} \\
\midrule
\multicolumn{3}{@{}l}{\footnotesize\textit{Essai court-circuit}} \\
$R_s = P_{cc}/I_{cc}^2$ & Res. ramenee & $\Omega$ \\
$\varepsilon_{cc} = V_{cc}/V_{1n}$ & Tension rel. & - \\
$\Delta U \approx \varepsilon_R\cos\varphi + \varepsilon_X\sin\varphi$ & Chute tens. & - \\
\midrule
\multicolumn{3}{@{}l}{\footnotesize\textit{Rendement}} \\
$\eta = P_2/(P_2 + P_{fer} + P_{cu})$ & Rendement & - \\
$P_{fer} = P_{cu}$ & Max. $\eta$ & - \\
$\eta_{typ} > 95\,\%$ & Valeur usuelle & - \\
\bottomrule
\end{tabular}
\end{multicols}
\normalsize
\end{tcolorbox}

% -----------------------------------------------------------------------
%  3. QUESTIONS TYPES
% -----------------------------------------------------------------------

\begin{tcolorbox}[
  ficheSection=colElec,
  title={3. Questions types - extraites des sujets reels},
]

% -- QT1 -------------------------------------------------------------

% Mini-procedure de resolution type copie concours
\begin{tcolorbox}[questiontype,title={Procedure type copie concours (5 etapes)}]
\begin{enumerate}[leftmargin=1.5em]
  \item \textbf{Qualifier le modele} : hypotheses, domaine de validite, grandeurs d'entree/sortie.
  \item \textbf{Ecrire les lois} : equations de base avec symboles et unites SI.
  \item \textbf{Isoler l'inconnue} : expression litterale, puis application numerique tracee.
  \item \textbf{Controler l'ordre de grandeur} : coherence physique et unite finale.
  \item \textbf{Conclure concours} : phrase courte orientee decision technique.
\end{enumerate}
\end{tcolorbox}

% Visuel de methode : chemin de resolution CAPET
\begin{center}
\begin{tikzpicture}[font=\small, node distance=7mm and 6mm, >=Stealth]
  \node[draw, rounded corners, fill=colElec!8, inner sep=4pt] (n1) {Notions};
  \node[draw, rounded corners, fill=colElec!8, right=of n1, inner sep=4pt] (n2) {Methode};
  \node[draw, rounded corners, fill=colElec!8, right=of n2, inner sep=4pt] (n3) {Application annale};
  \node[draw, rounded corners, fill=colElec!8, right=of n3, inner sep=4pt] (n4) {Verification jury};
  \draw[->, thick, colElec] (n1) -- (n2);
  \draw[->, thick, colElec] (n2) -- (n3);
  \draw[->, thick, colElec] (n3) -- (n4);
\end{tikzpicture}
\end{center}
\begin{tcolorbox}[questiontype,
  title={QT1 - Rapport de transformation - redresseur P3 vers 600\,V DC \hfill\normalfont\footnotesize E1 sujet 2026}
]
\textbf{enonce.}
Le reseau triphase fournit $230/400\,\mathrm{V}$. Un transformateur etoile-toile alimente un redresseur P3 (3 diodes, cathodes communes). La tension moyenne de sortie est $\bar{u}_s = \dfrac{3\sqrt{6}}{2\pi} \cdot V_{eff,simple}$. Le cahier des charges impose $\bar{u}_s = 600\,\mathrm{V}$. Calculer $V_{s,simple}$ puis le rapport de transformation $m$.

\tcbline

\textbf{Corrige.}

\textbf{etape 1 - tension simple secondaire necessaire.}
\[
  V_{s,simple} = \frac{600 \times 2\pi}{3\sqrt{6}} = \frac{600 \times 6{,}283}{7{,}348} \approx 513\,\mathrm{V}
\]

\textbf{etape 2 - rapport de transformation.}
\[
  m = \frac{V_{s,simple}}{V_{1,simple}} = \frac{513}{230} \approx \mathbf{2{,}23}
\]
\textit{Rappel : en couplage etoile-toile, le rapport $m$ s'applique indifferemment aux tensions simples ou composees.}
\end{tcolorbox}

\vspace{4pt}

% -- QT2 -------------------------------------------------------------

\begin{tcolorbox}[questiontype,
  title={QT2 - Rapport de transformation - redresseur PD6 vers 600\,V DC \hfill\normalfont\footnotesize E1 sujet 2026}
]
\textbf{enonce.}
Pour la solution PD6 (6 diodes), la tension moyenne est $\bar{u}_s = 0{,}95 \cdot U_M$ oe $U_M$ est l'amplitude des tensions composees secondaires. Pour obtenir $\bar{u}_s = 600\,\mathrm{V}$ avec un primaire $400\,\mathrm{V}$ composes, calculer $U_M$, puis $U_{s,composee}$, puis $m$.

\tcbline

\textbf{Corrige.}

\textbf{etape 1 - amplitude des tensions composees.}
\[
  U_M = \frac{600}{0{,}95} \approx 631\,\mathrm{V}
\]

\textbf{etape 2 - valeur efficace secondaire.}
\[
  U_{s,composee} = \frac{U_M}{\sqrt{2}} = \frac{631}{1{,}414} \approx 446\,\mathrm{V}
\]

\textbf{etape 3 - rapport de transformation.}
\[
  m = \frac{U_{s,composee}}{U_{1,composee}} = \frac{446}{400} \approx \mathbf{1{,}12}
\]
\textit{Comparaison : P3 necessite $m \approx 2{,}23$ contre $1{,}12$ pour PD6 - le PD6 exploite mieux le transformateur.}
\end{tcolorbox}

\vspace{4pt}

% -- QT3 -------------------------------------------------------------

\begin{tcolorbox}[questiontype,
  title={QT3 - Essai - vide - calcul de $R_f$ et $X_m$ \hfill\normalfont\footnotesize E2 BAC3 2026 (analogue MAS - $g\approx0$)}
]
\textbf{enonce.}
Lors de l'essai - vide d'un transformateur (ou d'une MAS - $g \approx 0$) : $P_0 = 7{,}82\,\mathrm{kW}$, $V_1 = 133\,\mathrm{V}$ (tension simple). Calculer $R_f$. Si $I_0 = 42\,\mathrm{A}$, calculer $Q_0$ puis $X_m$.

\tcbline

\textbf{Corrige.}

\textbf{etape 1 - resistance des pertes fer.}
\[
  R_f = \frac{V_1^2}{P_0} = \frac{133^2}{7820} = \frac{17689}{7820} \approx \mathbf{2{,}26\,\Omega}
\]

\textbf{etape 2 - puissance reactive.}
\[
  S_0 = V_1 \cdot I_0 = 133 \times 42 = 5586\,\mathrm{VA}
  \qquad
  Q_0 = \sqrt{S_0^2 - P_0^2} = \sqrt{5586^2 - 7820^2}
\]
\textit{Attention : verifier que $S_0 > P_0$ (sinon donnees incoherentes). Si $I_0$ n'est pas fourni, $Q_0$ ne peut etre calcule que si deux des trois grandeurs $P_0$, $Q_0$, $S_0$ sont donnees.}

\textbf{etape 3 - reactance magnetisante.}
\[
  X_m = \frac{V_1^2}{Q_0}
\]
\end{tcolorbox}

\end{tcolorbox}

% -----------------------------------------------------------------------
%  4. ATTENTES DU JURY
% -----------------------------------------------------------------------

\begin{tcolorbox}[
  jury,
  title={4. Attentes du jury - Erreurs recurrentes et points de vigilance},
]

\begin{enumerate}[label=\textbf{\arabic*.}, leftmargin=*, itemsep=5pt]

  \item \textbf{Confondre tensions simples et composees.}
  En couplage etoile-toile : $m = U_{s,comp}/U_{1,comp} = V_{s,simple}/V_{1,simple}$ - les deux sont equivalents.
  En couplage etoile-triangle, $m$ relie la tension simple primaire - la tension composee secondaire (ou vice-versa).

  \item \textbf{Oublier le facteur $\sqrt{3}$.}
  $U_{composee} = \sqrt{3} \cdot V_{simple}$ en etoile. Ne pas confondre les deux niveaux de tension.

  \item \textbf{Inverser essai - vide et essai court-circuit.}
  - vide - pertes \textbf{fer} ($R_f$, $X_m$), courant $I_0$ faible, tension nominale.
  Court-circuit - pertes \textbf{cuivre} ($R_s$), tension reduite $V_{cc}$, courant nominal.

  \item \textbf{Signe de $Q_0$.}
  $Q_0 = \sqrt{S_0^2 - P_0^2}$ est toujours \textbf{positif} (charge inductive).
  L'erreur classique est d'ecrire $Q_0 = \sqrt{P_0^2 - S_0^2}$, ce qui donne un resultat imaginaire.

  \item \textbf{Rapport $m$ et redresseur.}
  Le transformateur en amont d'un redresseur a un $m$ determine par la tension DC voulue.
  Calculer d'abord $U_{s,AC}$ necessaire depuis la formule $\bar{U}_{DC}$, puis $m = U_s/U_1$.

  \item \textbf{Justifier la demarche avant le calcul.}
  enoncer le modele, les hypotheses (frequence fixe, circuit magnetique lineaire), puis calculer.

\end{enumerate}

\medskip
\begin{tcolorbox}[
  colback=colJuryFrame!12!colJury,
  colframe=colJuryFrame,
  left=4pt, right=4pt, top=2pt, bottom=2pt,
  fontupper=\small\itshape
]
Rapport jury 2025 : "De nombreux candidats ne maetrisent pas les lois basiques de l'electricite (loi des neuds, loi des mailles, calculs de puissance ou encore la loi d'Ohm)."\\
Rapport jury 2024 : "Les candidats confondent souvent les couplages etoile/triangle et les rapports de transformation associes - verifier si $m$ porte sur les tensions simples ou composees selon le couplage."
\end{tcolorbox}

\end{tcolorbox}

% -----------------------------------------------------------------------
%  PIED - PReSENCE DANS LES SUJETS
% -----------------------------------------------------------------------

\begin{tcolorbox}[
  enhanced, breakable,
  colback=colGrayLight,
  colframe=colSeparator,
  left=4pt, right=4pt, top=3pt, bottom=3pt,
  before skip=4pt, after skip=0pt,
  fonttitle=\small\bfseries,
  title={Presence dans les sujets analyses (2013-2026)},
]
\begin{center}
\footnotesize\renewcommand{\arraystretch}{1.15}
\begin{tabular}{lccccccccccccc}
\toprule
\textbf{Annee} & 13 & 14 & 15 & 16 & 17 & 18 & 19 & 20 & 21 & 22 & 23 & 24 & 25 \\
\midrule
E1 (commune)   & $\checkmark$ & $\checkmark$ & - & $\checkmark$ & - & $\checkmark$ & - & - & $\checkmark$ & - & $\checkmark$ & - & $\checkmark$ \\
E2 discipl. IE & - & $\checkmark$ & $\checkmark$ & - & $\checkmark$ & - & $\checkmark$ & $\checkmark$ & - & $\checkmark$ & - & $\checkmark$ & - \\
\bottomrule
\end{tabular}
\end{center}
\smallskip
\footnotesize\textbf{Sous-themes frequents :}
rapport de transformation (10/13) - essai - vide $R_f$/$X_m$ (6/13) - couplages triphases Yy/Yd (5/13) - redresseur + transformateur amont (4/13) - rendement (3/13)

\medskip
\textbf{Ordres de grandeur (reperes rapides).}
\begin{itemize}[leftmargin=1.3em]
  \item Reseau BT FR : tension phase-neutre typique \SI{230}{\volt}, tension phase-phase typique \SI{400}{\volt}, frequence \SI{50}{\hertz}.
  \item Courants d'etude frequents en sujet : de quelques \si{\ampere} a quelques dizaines de \si{\ampere}.
  \item En conversion d'energie, un rendement attendu reste generalement entre 0{,}8 et 0{,}98 selon la technologie.
\end{itemize}

\medskip
\textbf{Synthese express (a refaire sans support).}
\begin{enumerate}[leftmargin=1.4em]
  \item Poser les hypotheses du modele et les conventions de signe.
  \item Ecrire 2 a 4 relations cles avec unites.
  \item Realiser un calcul numerique complet et verifier l'ordre de grandeur.
  \item Justifier le resultat par un argument technique de niveau CAPET.
\end{enumerate}

\end{tcolorbox}

\CAPETENDFILE









